\chapter{Reinforcement Learning for Execution and Market Making}\label{ch:ml:reinforcement-learning-for-execution-and-market-making}

Execution and market making both have optimal policies in closed form in their textbook models: the Almgren--Chriss
schedule (Book~10, chapter~14) and the inventory-aware quotes of Book~11 (chapter~3). The reason to learn a \omterm{def:ml:reinforcement-learning-foundations:mdp}{policy} instead
is everything the textbook model leaves out; the risk, as chapter~17 showed, is that a learned \omterm{def:ml:reinforcement-learning-foundations:mdp}{policy} exploits what its
training world gets wrong. This chapter measures both on problems whose optimum is known. A \omterm{def:ml:reinforcement-learning-for-execution-and-market-making:dqn}{deep Q-network} learns the
Almgren--Chriss schedule to within 0.4 basis points; moved to a market whose impact is three times higher, it loses 9.1
basis points to that market's optimum, four times what the textbook schedule of the wrong model loses. Trained on randomised
impact and allowed to see the impact of its own fills, it loses 0.9 there and 0.6 at home. A tabular learner for market
making, given 1.6 million steps, still falls short of the closed form and of a constant quote.

\section{State, action and reward design}

\begin{definition}[Reward shaping, action masking]\label{def:ml:reinforcement-learning-for-execution-and-market-making:design}
\emph{Reward shaping}\index{reward shaping} changes the \omterm{def:ml:reinforcement-learning-foundations:mdp}{rewards} an agent learns from without changing the \omterm{def:ml:reinforcement-learning-foundations:mdp}{policy} that is
optimal, for instance by removing a term whose expectation is known to be zero, or by adding a potential difference, to make
learning faster or more stable. \emph{Action masking}\index{action masking} removes the actions that are not allowed in a
state (selling more than is held, quoting a side that would breach an inventory limit) from the agent's choice and from the
maximum in its learning target.
\end{definition}

Design decides most of what an agent can learn. The execution agent sells one unit over ten steps in lots of a twentieth;
its state is the time, the holdings and one number described below; its action is the number of lots to sell, with every
sale above the holdings masked and the whole remainder forced at the last step. Its \omterm{def:ml:reinforcement-learning-foundations:mdp}{reward} is the Almgren--Chriss objective
split by step: minus the temporary-impact cost $\eta u^2/\tau$ of the sale, minus the risk charge $\lambda\sigma^2\tau x^2$ on
what is still held, with $\eta = 0.001$, $\sigma = 0.02$ and $\lambda = 10$ (in basis points, TWAP costs 21.4 and the optimum
18.85). The realised P\&L of the holdings over the step, $\sigma\sqrt\tau Z x$, has expectation zero; leaving it out is
\omterm{def:ml:reinforcement-learning-for-execution-and-market-making:design}{reward shaping}. Leaving it in is not harmless: with that noise in its \omterm{def:ml:reinforcement-learning-foundations:mdp}{reward}, the same network after the same 1\,500
episodes learned to sell the same amount every step, TWAP exactly, 2.55 basis points from the optimum.

\section{Deep Q-networks and their stabilisers}

\begin{definition}[Deep Q-network, experience replay, target network]\label{def:ml:reinforcement-learning-for-execution-and-market-making:dqn}
A \emph{deep Q-network}\index{deep Q-network} (DQN) is \omterm{def:ml:reinforcement-learning-foundations:td}{Q-learning} (chapter~17) with the \omterm{def:ml:reinforcement-learning-foundations:bellman}{action-value function} represented by
a \omterm{def:ml:neural-networks-for-noisy-tabular-data:mlp}{neural network}, trained by stochastic gradient steps on the temporal-difference error (Mnih and co-authors, 2015). Two
stabilisers make it work: \emph{experience replay}\index{experience replay} stores past transitions and trains on random
batches of them, breaking the correlation of consecutive steps; a \emph{target network}\index{target network} is a copy of the
network, updated only every few hundred steps, used to compute the learning targets so that they do not move with every
update.
\end{definition}

The chapter's DQN (\cref{lst:ml:rlt:dqn}) is a two-layer network of 64 units with 21 outputs, a replay memory of 20\,000
transitions, a \omterm{def:ml:reinforcement-learning-for-execution-and-market-making:dqn}{target network} copied every 200 updates and exploration decaying from 1 to 0.05 over the first half of
training. After 1\,500 episodes its greedy schedule costs 19.29 basis points, 0.44 above the continuous Almgren--Chriss
optimum; 0.20 of that gap is the lot grid, whose own optimum costs 19.05. Maximisation over noisy estimates biases \omterm{def:ml:reinforcement-learning-foundations:td}{Q-learning}
upwards; double \omterm{def:ml:reinforcement-learning-foundations:td}{Q-learning}, which chooses the action with one network and values it with the other (van Hasselt, Guez and
Silver, 2016), is the usual remedy and the design Ning and co-authors (2021) used for execution.

\section{Learning to execute}

\begin{definition}[Domain randomisation]\label{def:ml:reinforcement-learning-for-execution-and-market-making:dr}
\emph{Domain randomisation}\index{domain randomisation} trains a \omterm{def:ml:reinforcement-learning-foundations:mdp}{policy} on many versions of a simulator whose uncertain
parameters are drawn at random for each episode, so that the \omterm{def:ml:reinforcement-learning-foundations:mdp}{policy} works across them instead of exploiting one (Tobin and
co-authors, 2017).
\end{definition}

The textbook's weakness is its parameters: impact is estimated with error and changes with the market. Three worlds test the
agents: the model's ($\eta = 0.001$), a market with three times the impact, and one with a third of it; each has its own
Almgren--Chriss optimum, and every schedule is scored exactly against it (\cref{fig:ml:rlt:gaps}). The third state variable
is the impact the agent has observed in its own fills so far, as the logarithm of its ratio to the model's value (zero before
the first sale): a real execution algorithm measures its impact as it trades.

\begin{omfigure}
\begin{tikzpicture}
\begin{axis}[width=12cm,height=6cm,ybar,bar width=7pt,ymin=0,ymax=10,ylabel={gap to that world's optimum (bp)},
  area legend,xtick={0,1,2},xticklabels={model impact,impact $\times3$,impact $/3$},tick label style={font=\small},
  label style={font=\small},xmin=-0.5,xmax=2.5,grid=major,grid style={gray!25},legend columns=3,legend cell align=left,
  legend style={font=\footnotesize,at={(0.5,-0.18)},anchor=north,draw=none,fill=none},table/col sep=comma]
\addplot[fill=black!60,draw=none,bar shift=-16pt] table[x=i,y=acmodel]{figdata/ml/18-reinforcement-learning-for-execution-and-market-making/gaps.csv};
\addplot[fill=gray!40,draw=none,bar shift=-8pt] table[x=i,y=twap]{figdata/ml/18-reinforcement-learning-for-execution-and-market-making/gaps.csv};
\addplot[fill=omThm,draw=none,bar shift=0pt] table[x=i,y=dqn]{figdata/ml/18-reinforcement-learning-for-execution-and-market-making/gaps.csv};
\addplot[fill=omThm!40,draw=none,bar shift=8pt] table[x=i,y=noisy]{figdata/ml/18-reinforcement-learning-for-execution-and-market-making/gaps.csv};
\addplot[fill=omDef,draw=none,bar shift=16pt] table[x=i,y=dr]{figdata/ml/18-reinforcement-learning-for-execution-and-market-making/gaps.csv};
\legend{Almgren--Chriss (model),TWAP,DQN,DQN noisy reward,DQN randomised}
\end{axis}
\end{tikzpicture}

\omcaption{Each schedule's objective minus the Almgren--Chriss optimum of the world it trades in (basis points). The closed
form and the first DQN were built for the model's impact; the randomised DQN was trained on impact drawn over a factor of
three either way. Data: \texttt{ml\_rltrade.execution}.}\label{fig:ml:rlt:gaps}
\end{omfigure}

The model-trained DQN is the chapter's warning. In its own world it is 0.44 from the optimum; where impact is three times
higher it is 9.09 away, and where impact is lower 4.60: it has never seen its third input take those values, and a network
off its training distribution does something arbitrary. The closed form built on the same wrong impact loses 2.20 and 1.21:
a formula degrades smoothly where a network need not. The randomised DQN, trained for 3\,000 episodes with the impact drawn
anew each episode, learned to read its fills: it sells a little more slowly than the model's optimum at first and then
follows the impact it observes, and its gaps are 0.55, 0.86 and 1.24: the best of all in the high-impact world, and within 0.03 of the best (the
closed form, by luck of the direction of its error) in the low-impact one.
\Cref{fig:ml:rlt:schedules} shows the schedules in the high-impact world.

\begin{omfigure}
\begin{tikzpicture}
\begin{axis}[width=11cm,height=6cm,xlabel={step},ylabel={holdings (share of the order)},tick label style={font=\small},
  label style={font=\small},xmin=0,xmax=10,ymin=0,ymax=1.02,grid=major,grid style={gray!25},legend columns=3,legend cell align=left,
  legend style={font=\footnotesize,at={(0.5,-0.25)},anchor=north,draw=none,fill=none,
  /tikz/every even column/.append style={column sep=6pt}},table/col sep=comma]
\addplot[black,ultra thick] table[x=step,y=optimum]{figdata/ml/18-reinforcement-learning-for-execution-and-market-making/schedules.csv};
\addplot[black!50,thick,dashed] table[x=step,y=acmodel]{figdata/ml/18-reinforcement-learning-for-execution-and-market-making/schedules.csv};
\addplot[gray,thick,dotted] table[x=step,y=twap]{figdata/ml/18-reinforcement-learning-for-execution-and-market-making/schedules.csv};
\addplot[omThm,very thick,mark=*,mark size=1.4pt] table[x=step,y=dqn]{figdata/ml/18-reinforcement-learning-for-execution-and-market-making/schedules.csv};
\addplot[omDef,very thick,mark=square*,mark size=1.4pt] table[x=step,y=dr]{figdata/ml/18-reinforcement-learning-for-execution-and-market-making/schedules.csv};
\legend{optimum for impact $\times3$,Almgren--Chriss (model),TWAP,DQN,DQN randomised}
\end{axis}
\end{tikzpicture}

\omcaption{Holdings over the ten steps in the world with three times the model's impact. Data:
\texttt{ml\_rltrade.agents}.}\label{fig:ml:rlt:schedules}
\end{omfigure}

\section{Learning to make markets}

Market making in the model of Book~11 (chapter~3) has an exact solution: a mid-price that diffuses, fills at depth $\delta$
with intensity $Ae^{-k\delta}$, a running penalty $\phi q^2$ on inventory and a terminal penalty $\alpha q^2$, here with
$A = 140$, $k = 1.5$, $\sigma = 2$, $\phi = 0.02$, $\alpha = 0.01$ and inventory within $\pm10$. The optimal depths at the start
are 0.68 on both sides when flat and 0.57 and 0.79 when five units long (tighter on the side that reduces inventory). The
learner is tabular \omterm{def:ml:reinforcement-learning-foundations:td}{Q-learning} (\cref{lst:ml:rlt:mm}) over 20 time buckets and the 21 inventories, choosing each side's depth
among six values, with the zero-mean inventory P\&L shaped out of its \omterm{def:ml:reinforcement-learning-foundations:mdp}{reward} and inventory-breaching sides masked.
\Cref{tab:ml:rlt:mm} scores every \omterm{def:ml:reinforcement-learning-foundations:mdp}{policy} with the firm's market-making simulator on 5\,000 common paths.

\begin{table}[htbp]
\centering\small
\begin{tabular}{@{}lccccc@{}}
\toprule
\omterm{def:ml:reinforcement-learning-foundations:mdp}{policy} & objective & (s.e.) & P\&L s.d. & $|q_T|$ & fills\\
\midrule
exact optimum (Book~11) & 67.95 & 0.12 & 8.53 & 2.67 & 101.8\\
Avellaneda--Stoikov, $\gamma = 0.1$ & 64.73 & 0.09 & 6.50 & 2.26 & 97.0\\
constant depth $1/k$ & 66.67 & 0.17 & 11.61 & 5.03 & 101.0\\
\omterm{def:ml:reinforcement-learning-foundations:td}{Q-learning}, 2\,000 episodes & 62.34 & 0.12 & 8.31 & 4.40 & 91.7\\
\omterm{def:ml:reinforcement-learning-foundations:td}{Q-learning}, 8\,000 episodes & 62.84 & 0.12 & 8.53 & 4.38 & 93.1\\
\bottomrule
\end{tabular}
\caption{Market making for one unit of time: mean of P\&L minus the running and terminal inventory penalties (with its
standard error), P\&L standard deviation, mean closing inventory and fills per path, on 5\,000 paths with common random
numbers. Data: \texttt{ml\_rltrade.market\_making}.}\label{tab:ml:rlt:mm}
\end{table}

The learner improves slowly and stays below the constant quote: 1.6 million steps of experience are not enough to rank six
depths whose expected values differ by less than the noise of a fill. The spread earned per step is a small signal under a
large noise, which is the market maker's problem in general. Closed forms or parametric policies fitted by \omterm{def:ml:reinforcement-learning-foundations:mdp}{policy} search in
a simulator (a depth that is a line in inventory has two parameters, where the table has 420 states and 36 actions) are what work at this scale;
learning earns its place on the parts the model leaves out, such as queue position, adverse selection and several venues,
where no closed form exists.

\section{What works in the published evidence}

The evidence is thinner than the attention. Nevmyvaka, Feng and Kearns (2006) applied \omterm{def:ml:reinforcement-learning-foundations:mdp}{reinforcement learning} to execution on
a year and a half of millisecond limit-order data from NASDAQ, with a state space factorised to keep learning tractable;
Ning and co-authors (2021) trained a double \omterm{def:ml:reinforcement-learning-for-execution-and-market-making:dqn}{deep Q-network} with \omterm{def:ml:reinforcement-learning-for-execution-and-market-making:dqn}{experience replay} on order-book features for nine stocks and
found it beat the standard benchmark on most of them; Spooner and co-authors
(2018) built a high-fidelity limit-order-book simulator and a temporal-difference market maker with a custom \omterm{def:ml:reinforcement-learning-foundations:mdp}{reward} that
controls inventory risk, which beat simple benchmarks and an online-learning approach in that simulator. Hambly, Xu and Yang
(2023) survey the field. The common thread is the simulator: results are reported in a simulator, often the authors' own, and
the transfer to live trading is rarely documented in public. A desk should read them as evidence that the methods work where
the simulator is right, and build its own evidence with the tests of chapter~17 and this chapter: exact benchmarks, several
worlds, \omterm{def:ml:reinforcement-learning-foundations:ope}{off-policy evaluation} on its own logs.

\begin{method}[Deploying a learned execution or quoting policy]\label{met:ml:rlt:method}
\begin{enumerate}
\item Benchmark against the closed form of the textbook model in that model; a learner that cannot match it there is not
ready.
\item Shape \omterm{def:ml:reinforcement-learning-foundations:mdp}{rewards} to remove known zero-mean noise, mask forbidden actions, and give the state the signals a trader would
use (fills, observed impact, queue position).
\item Train on randomised simulators whose parameters span their estimation error, and evaluate in held-out parameter
settings, never only in the training one.
\item Bound the \omterm{def:ml:reinforcement-learning-foundations:mdp}{policy} (participation limits, inventory limits, a fallback to the closed form) and compare live against the
incumbent by \omterm{def:ml:reinforcement-learning-foundations:ope}{off-policy evaluation} and small controlled trials.
\end{enumerate}
\end{method}

\section{Tutorial: beating a closed form}\label{tut:ml:reinforcement-learning-for-execution-and-market-making:tut}

\begin{tutorial}
\textbf{Goal.} Train a DQN on the Almgren--Chriss liquidation with and without \omterm{def:ml:reinforcement-learning-for-execution-and-market-making:dr}{domain randomisation} and score it in three
worlds; train a tabular market maker and score it against the exact optimum. \textbf{End state:}
\cref{fig:ml:rlt:gaps,fig:ml:rlt:schedules}, \cref{tab:ml:rlt:mm}.
\begin{tutsteps}
\item \textbf{A DQN update with replay, a \omterm{def:ml:reinforcement-learning-for-execution-and-market-making:dqn}{target network} and masking.}
\omcode{code/firm/rltrade/firm_rltrade.py}{113}{143}{The training loop of deep Q-learning.}\label{lst:ml:rlt:dqn}
\item \textbf{\omterm{def:ml:reinforcement-learning-foundations:td}{Q-learning} for market making.}
\omcode{code/firm/rltrade/firm_rltrade.py}{228}{243}{Tabular Q-learning with decaying step sizes.}\label{lst:ml:rlt:mm}
\item \textbf{Run} \texttt{ml\_rltrade.execution()} (about a minute on one core), \texttt{market\_making()} and
\texttt{fig\_rltrade.py}.
\end{tutsteps}
\textbf{What to change next.} Use double \omterm{def:ml:reinforcement-learning-foundations:td}{Q-learning}; replace the market maker's table by a two-parameter \omterm{def:ml:reinforcement-learning-foundations:mdp}{policy} fitted by
\omterm{def:ml:reinforcement-learning-foundations:mdp}{policy} search; add a passive-order action whose fills come from queue position in Book~10's exchange simulator.
\end{tutorial}

\section{Build: trading environments and learners}\label{bld:ml:reinforcement-learning-for-execution-and-market-making:rltrade}

\begin{build}
\bfield{Purpose} Execution and market-making policies learned where their optimum is known, and tested away from it.
\bfield{Interface} \texttt{ACEnv(X, T, n, lots, eta, sigma, lam, eta\_range, noise)} with \texttt{reset}, \texttt{step},
\texttt{mask}; \texttt{DQN}, \texttt{train\_dqn(env, episodes, seed)}, \texttt{dqn\_schedule(env, net, eta)},
\texttt{grid\_optimum}, \texttt{ac\_objective} (on \texttt{firm.acexec}); \texttt{MMEnv}, \texttt{q\_learning\_mm},
\texttt{mm\_policy}, \texttt{mm\_objective} (with \texttt{firm.invmm.simulate}).
\bfield{Rules} Every learned \omterm{def:ml:reinforcement-learning-foundations:mdp}{policy} is scored exactly or on common random numbers against the closed form of the world it
is tested in; training parameters and test parameters are reported separately.
\bfield{Acceptance tests} \texttt{code/firm/rltrade/tests/}: the environment's \omterm{def:ml:reinforcement-learning-foundations:mdp}{rewards} sum to the Almgren--Chriss objective of
the schedule followed; masking never allows an illegal sale; the lot-grid optimum is no worse than any enumerated schedule;
DQN training is deterministic and beats TWAP in the model world; the market-making environment's inventory changes match
fills with probability $Ae^{-k\delta}\,dt$ per side.
\bfield{Stretch} Double DQN; \omterm{def:ml:reinforcement-learning-foundations:mdp}{policy} search for quoting; the exchange-simulator wrapper (Book~10, chapter~26) with queue
position as state.
\end{build}

\begin{omsources}
\item V.~Mnih and co-authors, ``Human-level control through deep reinforcement learning'', \emph{Nature} 518, 2015.
\item H.~van Hasselt, A.~Guez and D.~Silver, ``Deep reinforcement learning with double Q-learning'', arXiv:1509.06461, 2016.
\item J.~Tobin and co-authors, ``Domain randomization for transferring deep neural networks from simulation to the real
world'', arXiv:1703.06907, 2017.
\item Y.~Nevmyvaka, Y.~Feng and M.~Kearns, ``Reinforcement learning for optimized trade execution'', ICML, 2006.
\item B.~Ning and co-authors, ``Double deep Q-learning for optimal execution'', \emph{Applied Mathematical Finance}, 2021.
\item T.~Spooner, J.~Fearnley, R.~Savani and A.~Koukorinis, ``Market making via reinforcement learning'', arXiv:1804.04216,
2018.
\item B.~Hambly, R.~Xu and H.~Yang, ``Recent advances in reinforcement learning in finance'', arXiv:2112.04553, 2023.
\item R.~Almgren and N.~Chriss, ``Optimal execution of portfolio transactions'', \emph{Journal of Risk} 3(2), 2001.
\item M.~Avellaneda and S.~Stoikov, ``High-frequency trading in a limit order book'', \emph{Quantitative Finance} 8(3), 2008.
\end{omsources}

\section{Exercises}

\begin{exercise}[$\star$]\label{exo:ml:reinforcement-learning-for-execution-and-market-making:1}
Compute TWAP's objective in the model world: ten sales of 0.1 with $\eta = 0.001$, $\tau = 0.1$, and the risk charge with
$\sigma = 0.02$, $\lambda = 10$. Give it in basis points.
\end{exercise}

\begin{exercise}[$\star$]\label{exo:ml:reinforcement-learning-for-execution-and-market-making:2}
At what depth does $\delta Ae^{-k\delta}$, the expected spread earned per unit time on one side, peak? Compare with the exact
optimum's flat depth.
\end{exercise}

\begin{exercise}[$\star$]\label{exo:ml:reinforcement-learning-for-execution-and-market-making:3}
Why must the mask also enter the target $\max_{u'}Q(x', u')$, not only the choice of action?
\end{exercise}

\begin{exercise}[$\star\star$]\label{exo:ml:reinforcement-learning-for-execution-and-market-making:4}
Show that removing $\sigma\sqrt\tau Zx$ from each step's \omterm{def:ml:reinforcement-learning-foundations:mdp}{reward} leaves the optimal \omterm{def:ml:reinforcement-learning-foundations:mdp}{policy} unchanged. Why did the noisy agent
end at TWAP?
\end{exercise}

\begin{exercise}[$\star\star$]\label{exo:ml:reinforcement-learning-for-execution-and-market-making:5}
Why did the model-trained DQN do worse than the wrong closed form in the high-impact world?
\end{exercise}

\begin{exercise}[$\star\star$]\label{exo:ml:reinforcement-learning-for-execution-and-market-making:6}
\emph{Find the flaw.} ``We trained the market maker with \omterm{def:ml:reinforcement-learning-for-execution-and-market-making:dr}{domain randomisation} over volatility and it was profitable in every
one of 10\,000 randomised simulator runs, so it is robust.''
\end{exercise}

\begin{exercise}[$\star\star\star$]\label{exo:ml:reinforcement-learning-for-execution-and-market-making:7}
\emph{Coding.} Replace the market maker's table by the \omterm{def:ml:reinforcement-learning-foundations:mdp}{policy} $\delta_{b,a} = c\pm dq$ and choose $(c, d)$ on a grid by the
simulator's mean objective on 2\,000 paths. How close do you get to the exact optimum?
\end{exercise}

\begin{exercise}[$\star\star\star$]\label{exo:ml:reinforcement-learning-for-execution-and-market-making:8}
Derive the discrete Almgren--Chriss schedule for the chapter's objective by writing the first-order conditions in the
holdings $x_1,\dots,x_{n-1}$.
\end{exercise}

\section{Problem: Beating a Closed Form}

\begin{problem}[{Weekend problem --- when to learn a policy}]\label{pb:ml:reinforcement-learning-for-execution-and-market-making:1}
The chapter's liquidation and market-making problems.

\medskip\noindent\textbf{Part I --- Design.}
\begin{enumerate}
\item What are the execution agent's state, actions and \omterm{def:ml:reinforcement-learning-foundations:mdp}{reward}?
\item Why is the P\&L noise removed from the \omterm{def:ml:reinforcement-learning-foundations:mdp}{reward}, and what happened when it was not?
\item What does \omterm{def:ml:reinforcement-learning-for-execution-and-market-making:design}{action masking} prevent?
\item Why give the agent the impact it has observed?
\end{enumerate}

\medskip\noindent\textbf{Part II --- Execution.}
\begin{enumerate}[resume]
\item What do TWAP, the lot-grid optimum and the DQN cost in the model world?
\item What are the gaps of each schedule in the three worlds?
\item Why does the randomised agent lose a little at home?
\item What would change with a price that trends during the order?
\end{enumerate}

\medskip\noindent\textbf{Part III --- Market making.}
\begin{enumerate}[resume]
\item What are the exact optimum's depths at the start?
\item How do the five policies score?
\item Why does the tabular learner stay below a constant quote?
\item What would you learn, and what would you keep in closed form?
\end{enumerate}

\medskip\noindent\textbf{Part IV --- The verdict.}
\begin{enumerate}[resume]
\item State the \emph{named result}: the learned \omterm{def:ml:reinforcement-learning-foundations:mdp}{policy}'s shortfall against the Almgren--Chriss optimum in the model world and
in the mis-specified worlds, with and without \omterm{def:ml:reinforcement-learning-for-execution-and-market-making:dr}{domain randomisation}.
\item What does the published evidence show, and what does it not?
\item How would you deploy the randomised execution agent?
\item What is the fallback if the agent's observed impact leaves its training range?
\item Why does a table not scale to a real order book?
\item Which stabiliser matters most in the DQN here, and how would you check?
\item What should the comparison with the incumbent algorithm look like?
\item In one sentence: when is a learned \omterm{def:ml:reinforcement-learning-foundations:mdp}{policy} worth more than a closed form?
\end{enumerate}
\end{problem}

\section{Interview questions}

\begin{interviewq}[$\star$ \iqroles{mle}]\label{iq:ml:reinforcement-learning-for-execution-and-market-making:1}
What are \omterm{def:ml:reinforcement-learning-for-execution-and-market-making:dqn}{experience replay} and \omterm{def:ml:reinforcement-learning-for-execution-and-market-making:dqn}{target networks} for?
\end{interviewq}

\begin{interviewq}[$\star\star$ \iqroles{researcher, trader}]\label{iq:ml:reinforcement-learning-for-execution-and-market-making:2}
How would you frame optimal execution as a reinforcement-learning problem? What goes in the state?
\end{interviewq}

\begin{interviewq}[$\star\star$ \iqroles{researcher}]\label{iq:ml:reinforcement-learning-for-execution-and-market-making:3}
What is \omterm{def:ml:reinforcement-learning-for-execution-and-market-making:design}{reward shaping}, and how can it go wrong?
\end{interviewq}

\begin{interviewq}[$\star\star$ \iqroles{mle}]\label{iq:ml:reinforcement-learning-for-execution-and-market-making:4}
Your DQN's Q-values keep growing during training. What is happening and what do you do?
\end{interviewq}

\begin{interviewq}[$\star\star$ \iqroles{researcher, trader}]\label{iq:ml:reinforcement-learning-for-execution-and-market-making:5}
Would you use \omterm{def:ml:reinforcement-learning-foundations:mdp}{reinforcement learning} for market making? Where would it help and where not?
\end{interviewq}

\begin{interviewq}[$\star\star\star$ \iqroles{researcher}]\label{iq:ml:reinforcement-learning-for-execution-and-market-making:6}
What is \omterm{def:ml:reinforcement-learning-for-execution-and-market-making:dr}{domain randomisation}, and why can it help and hurt a trading \omterm{def:ml:reinforcement-learning-foundations:mdp}{policy}?
\end{interviewq}
